\chapter{GitHub}\label{ch:git-github}

Git is the version control system; GitHub is a website that hosts Git
repositories and adds a collaboration layer on top --- pull requests, issues,
code review, continuous integration. The distinction matters: everything in
this chapter is GitHub's invention, not Git's, and GitLab, Gitea and Bitbucket
do the same things with different names.

\section{Repositories, Forks and Clones}\label{sec:gh-repos}

\begin{keys}[0.22]
	\vocab{Repository} & A project on GitHub: the Git repository plus issues, pull requests and settings \\
	\vocab{Clone}      & A local copy on your machine. Pure Git                        \\
	\vocab{Fork}       & \emph{Your} copy of someone else's repository, on GitHub. A GitHub concept \\
\end{keys}

\begin{intuition}
	Fork when you cannot push to the original --- which is the normal situation
	for contributing to a project you do not maintain. You fork on GitHub, clone
	your fork, push branches to it, and open a pull request asking the original to
	take your changes. When you \emph{can} push, skip the fork and branch directly.
\end{intuition}

\begin{shell}
$ git clone git@github.com:Konyi0613/LaTeX.git
$ cd LaTeX
$ git remote add upstream git@github.com:original/LaTeX.git   # for a fork
$ git remote -v
origin      git@github.com:Konyi0613/LaTeX.git (fetch)
upstream    git@github.com:original/LaTeX.git   (fetch)
\end{shell}

\section{Authentication}\label{sec:gh-auth}

GitHub stopped accepting passwords for Git operations in 2021. There are three
replacements.

\begin{keys}[0.26]
	\vocab{SSH key}     & A key pair; the public half goes to GitHub. No expiry, no typing \\
	\vocab{Personal access token} & A password-shaped secret with scopes and an expiry date, used over HTTPS \\
	\vocab{\cmd{gh} CLI} & Logs in once and acts as a credential helper for everything else \\
\end{keys}

\begin{shell}
$ ssh-keygen -t ed25519 -C "you@example.com"    # see sec. 15.2
$ cat ~/.ssh/id_ed25519.pub                     # paste into GitHub settings
$ ssh -T git@github.com                         # check it worked
Hi Konyi0613! You've successfully authenticated.
\end{shell}

\begin{note}[This machine]
	Authentication here is already handled by the third route: \cmd{gh} is
	registered as Git's credential helper for \texttt{github.com}, so HTTPS pushes
	find a token without prompting. \cmd{gh auth status} shows the account and
	scopes. An SSH key is still worth having --- it works when \cmd{gh} is not
	installed, which is every server you will ever log into.
\end{note}

\begin{moral}
	Never put a token in a URL, a script, or a commit. GitHub scans public pushes
	for its own token formats and revokes what it finds, which is a kindness rather
	than a solution --- \autoref{sec:git-bigfiles} covers what to do when a secret
	is already in the history.
\end{moral}

\section{Pull Requests}\label{sec:gh-pr}

\begin{definition}[Pull request]
	A request to merge one branch into another, with a description, a discussion,
	a running diff and whatever automated checks the project has configured. It is
	a GitHub object, not a Git one --- underneath it is just two branches.
\end{definition}

The full cycle:

\begin{shell}
$ git switch -c fix-config-crash        # 1. branch
$ # ... edit, commit ...
$ git push -u origin fix-config-crash   # 2. push the branch
$ gh pr create --fill                   # 3. open the PR
$ # ... review happens, you push more commits ...
$ gh pr merge --squash --delete-branch  # 4. merge and tidy up
\end{shell}

\begin{keys}[0.26]
	\vocab{Merge commit}  & Keeps every commit and records the branch. Full history  \\
	\vocab{Squash merge}  & Combines the branch into one commit on \cmd{main}. Tidy history, loses the detail \\
	\vocab{Rebase merge}  & Replays each commit onto \cmd{main}. Linear history, no merge commit \\
\end{keys}

\begin{note}
	Squash merging pairs well with a habit of committing untidily while you work:
	the mess stays in the pull request, where it is reviewable, and \cmd{main} gets
	one clean commit per change. The cost is that \cmd{git log} on \cmd{main} can
	no longer show you how the work was built up. Most projects pick one of the
	three and apply it consistently; find out which before your first PR.
\end{note}

\begin{note}[Draft pull requests]
	Open a PR as a draft as soon as you push the branch, not when the work is
	finished. It gives the change a URL, runs CI on every push, and tells your
	colleagues what you are doing --- with no implication that it is ready.
\end{note}

\section{Reviewing Code}\label{sec:gh-review}

\begin{keys}[0.26]
	\textbf{Comment}         & Feedback without a verdict                            \\
	\textbf{Approve}         & Fine to merge                                         \\
	\textbf{Request changes} & Blocks merging until resolved                         \\
	\textbf{Suggestion}      & A concrete edit the author can apply with one click    \\
\end{keys}

\begin{keys}[0.34]
	\cmd{gh pr list}                & Open pull requests                               \\
	\cmd{gh pr view 42}             & One, in the terminal                             \\
	\cmd{gh pr view 42 -{}-web}     & \dots{} in the browser                           \\
	\cmd{gh pr diff 42}             & Its diff                                         \\
	\cmd{gh pr checkout 42}         & Check the branch out locally --- run it, don't just read it \\
	\cmd{gh pr review 42 -{}-approve} & Approve from the command line                  \\
	\cmd{gh pr checks 42}           & Whether CI passed                                \\
\end{keys}

\begin{moral}
	\cmd{gh pr checkout} is the command that makes reviews useful. Reading a diff
	catches style and obvious mistakes; checking the branch out, building it, and
	running it catches the things that matter. A review that never left the browser
	has only checked what the diff shows.
\end{moral}

\begin{note}
	Review the \emph{change}, not the person, and say which comments are blocking
	and which are preferences. ``Consider extracting this'' and ``this will
	deadlock'' need to be distinguishable, or the author has to guess.
\end{note}

\section{Issues and Projects}\label{sec:gh-issues}

\begin{keys}[0.34]
	\cmd{gh issue list}                    & Open issues                              \\
	\cmd{gh issue create -t "..." -b "..."} & Create one                              \\
	\cmd{gh issue view 17}                 & Read one                                 \\
	\cmd{gh issue close 17}                & Close it                                 \\
	\cmd{gh issue list -{}-label bug -{}-assignee @me} & Filter                       \\
\end{keys}

\begin{note}[Closing keywords]
	A commit message or pull request body containing \texttt{Closes \#42},
	\texttt{Fixes \#42} or \texttt{Resolves \#42} closes that issue automatically
	when the change lands on the default branch --- and links the two permanently,
	so anyone reading the issue later finds the commit that fixed it. This is worth
	more than the automation: it is free documentation.
\end{note}

\section{The \texorpdfstring{\cmd{gh}}{gh} Command-Line Tool}\label{sec:gh-cli}

\begin{keys}[0.34]
	\cmd{gh auth login}, \cmd{gh auth status} & Sign in; check who you are          \\
	\cmd{gh repo clone owner/name}      & Clone without typing a URL                 \\
	\cmd{gh repo create -{}-public}     & Create a repository from here              \\
	\cmd{gh repo view -{}-web}          & Open it in a browser                       \\
	\cmd{gh pr create -{}-fill}         & Open a PR using the commits as the description \\
	\cmd{gh pr status}                  & Your PRs, and those awaiting your review   \\
	\cmd{gh run list}, \cmd{gh run watch} & Actions runs; follow one live            \\
	\cmd{gh release create v1.0}        & Cut a release                              \\
	\cmd{gh browse}                     & Open the current file or line on GitHub    \\
\end{keys}

\begin{note}
	\cmd{gh} 2.93 is installed here. \cmd{gh browse} deserves particular mention:
	run from a repository it opens the project page, and \cmd{gh browse
		Chapters/git\_1\_model.tex} opens that file --- which is the fastest way to get
	a shareable link to a specific line of code.
\end{note}

\section{GitHub Actions}\label{sec:gh-actions}

\begin{definition}[Actions]
	GitHub's continuous integration. A workflow is a YAML file, stored in the
	repository under \file{.github/workflows}, describing jobs to run on a hosted
	machine when something happens: a push, a pull request, a schedule, or a
	manual trigger.
\end{definition}

\begin{conf}
# .github/workflows/build.yml
name: Build the note

on:
  push:
    branches: [main]
  pull_request:

jobs:
  build:
    runs-on: ubuntu-latest
    steps:
      - uses: actions/checkout@v4

      - name: Install TeX Live
        run: |
          sudo apt-get update
          sudo apt-get install -y texlive-xetex texlive-latex-extra latexmk

      - name: Compile
        run: latexmk -xelatex -interaction=nonstopmode -output-directory=build CSTools.tex

      - name: Upload the PDF
        uses: actions/upload-artifact@v4
        with:
          name: CSTools.pdf
          path: build/CSTools.pdf
\end{conf}

\begin{keys}[0.24]
	\texttt{on}        & What triggers the workflow                                 \\
	\texttt{jobs}      & Independent units; they run in parallel by default          \\
	\texttt{runs-on}   & Which hosted machine                                        \\
	\texttt{steps}     & Run in order; each is a \texttt{run} command or a \texttt{uses} action \\
	\texttt{uses}      & A reusable action from the marketplace --- pin the version  \\
\end{keys}

\begin{moral}
	The first workflow worth writing for any project is the one that proves it
	still builds. For this note that is one \cmd{latexmk} invocation, and it means
	a broken \LaTeX{} file is caught by the pull request rather than by you, three
	commits later.
\end{moral}

\section{Pages, Releases and \texorpdfstring{\file{README}}{README} Conventions}\label{sec:gh-pages}

\begin{keys}[0.26]
	\file{README.md}        & Shown on the repository's front page. What this is, how to build it, how to use it \\
	\file{LICENSE}          & Without one, nobody may legally reuse your work        \\
	\file{CONTRIBUTING.md}  & How to submit a change                                 \\
	\file{.github/}         & Workflows, issue templates, PR templates               \\
	\vocab{Releases}        & Tagged versions with notes and attached binaries       \\
	\vocab{Pages}           & Static hosting from a branch or an Actions job         \\
\end{keys}

\begin{note}
	The \file{README} is the only documentation most people will ever read.
	Four things earn their space: what the project is in one sentence, how to
	install or build it, a minimal working example, and where to ask questions.
	Badges are decoration; the build instructions are not.
\end{note}
